% !TeX root = ../technical_report.tex
\chapter{Class-Balanced Loss Formulation and Long-Tail Gradient Dynamics}
\label{ch:class_balanced_loss}

% CHAPTER BRIEF FOR JULES / TECHNICAL WRITER:
% -------------------------------------------------------------
% Focus: Cui et al. (CVPR 2019) loss implementation and mathematical dynamics on Pareto distributions.
% Repository Files:
%   - src/taxon_vision/models/loss.py
% Mathematical Derivations & Code Details:
%   - Extreme Long-Tail Distribution: Zipfian / Pareto scaling of biological species sightings.
%   - Effective Number of Samples (Cui et al., 2019):
%       * E_n = (1 - \beta^n) / (1 - \beta) where \beta \in [0, 1) captures feature space volume overlap.
%       * As n -> \infty, E_n saturates at 1 / (1 - \beta).
%       * As n -> 1, E_n \approx n.
%   - Class-Balanced Weights:
%       * W_i = (1 - \beta) / (1 - \beta^{N_i})
%       * Normalisation: \sum_{i=1}^K W_i = K (mean class weight equals 1.0).
%   - Numerical Safety: clamping minimum sample count to 1.0 to avoid zero-division; PyTorch buffer registration.
%   - Loss Computation: F.cross_entropy(logits, targets, weight=device_weights).
%   - Contrast Analysis: why CB-CE beats unweighted cross-entropy (prevents tail representation collapse)
%     and naive inverse frequency (avoids explosive gradient spikes from rare noisy singletons).
% -------------------------------------------------------------

\section{The Biological Long-Tail Pareto Distribution}
\label{sec:biological_long_tail}

\section{Effective Number of Samples Mathematical Derivation}
\label{sec:effective_number_mathematics}

\section{PyTorch Implementation and Numerical Stabilisation}
\label{sec:cb_pytorch_implementation}
